\chapter{Introduction générale}
\label{ch:intro}

\section{Contexte général}

Le LIAS --- Laboratoire d'Informatique et Automatique pour les Systèmes --- est une
structure de recherche rattachée à la Faculté des Sciences Ben M'Sik de l'Université
Hassan~II de Casablanca. Comme toute unité de recherche, il fonctionne sur deux plans
distincts et rarement outillés ensemble.

Sur le plan externe, le laboratoire doit être visible : un visiteur, un candidat à une
thèse ou un partenaire industriel doit pouvoir consulter la composition des équipes, la
production scientifique et le calendrier des manifestations sans disposer d'un compte.

Sur le plan interne, le laboratoire produit en permanence de l'information administrative
dont la valeur est cumulative : des membres entrent, changent de statut, prennent des
responsabilités, partent puis reviennent parfois ; des mandats de direction se succèdent ;
des documents, des procès-verbaux et des conventions s'accumulent et doivent rester
retrouvables plusieurs années après. Cette information est, dans la pratique courante,
dispersée entre tableurs, messageries et dossiers partagés --- des supports qui ne
conservent ni l'ordre chronologique des décisions ni leur auteur.

\section{Problématique}

Le cahier des charges du projet formule une exigence qui dépasse la simple numérisation
de ces pratiques. Sa section~21 énonce quatre règles générales : aucun membre ne perd ses
données, tout est historisé, toute action importante est traçable, et surtout une
séparation stricte doit être maintenue entre le \srt{} d'appartenance d'un membre
(permanent, associé, doctorant\dots) et son \rol{} de responsabilité.

Cette dernière règle est la véritable difficulté du sujet. Dans une application de gestion
classique, le rôle applicatif est un champ que l'administrateur renseigne. Ici, il ne peut
pas l'être, pour une raison métier : un enseignant-chercheur \emph{permanent} qui devient
directeur du laboratoire ne cesse pas d'être permanent ; il acquiert, pour une période
datée, une responsabilité supplémentaire. À la fin de son mandat, il redevient un membre
permanent ordinaire sans que rien n'ait été effacé. Trois notions coexistent donc et
évoluent à des rythmes différents :

\begin{itemize}
  \item le \srt{} --- une catégorie d'appartenance, qui change lentement ;
  \item le \man{} --- une responsabilité de gouvernance bornée par deux dates ;
  \item le \rol{} --- l'ensemble des permissions que la sécurité applicative lit à chaque
        requête.
\end{itemize}

La problématique du projet se formule alors ainsi : \emph{comment construire un système
d'information dans lequel les permissions techniques découlent automatiquement et de
manière vérifiable de la réalité institutionnelle du laboratoire, sans jamais détruire
l'historique qui les justifie ?}

\section{Objectifs}

Les objectifs poursuivis découlent directement de la section~1 du cahier des charges et
sont repris au tableau~\ref{tab:objectifs}.

\begin{table}[H]
\centering
\caption{Objectifs du système et traduction technique retenue}
\label{tab:objectifs}
\small
\begin{tabularx}{\textwidth}{@{}p{5.1cm}X@{}}
\toprule
\textbf{Objectif (CdC §1)} & \textbf{Traduction dans le système} \\
\midrule
Centraliser la gestion du laboratoire & Une base relationnelle unique, 22 contrôleurs REST, une interface unifiée pour 18 modules métier \\
Assurer la traçabilité des activités & Journal d'audit systématique (acteur, action, entité, horodatage, adresse IP) \\
Capitaliser la mémoire du laboratoire & Suppression logique généralisée, quatre tables d'historique, archivage documentaire versionné \\
Structurer membres et responsabilités & Séparation \srt{}~/~\rol{}~/~\man{} avec dérivation automatique du rôle \\
Faciliter le rapport annuel obligatoire & Agrégation automatique des indicateurs et export PDF, déclenché manuellement ou par planification annuelle \\
Améliorer la collaboration & Discussions à trois portées, notifications ciblées, calendrier unifié \\
\bottomrule
\end{tabularx}
\end{table}

\section{Présentation du projet LIAS}

La plateforme réalisée est une application web unique servant deux publics. Sa partie
publique, accessible sans authentification, présente le laboratoire, ses équipes, ses
publications et ses événements, et permet le dépôt d'une candidature d'adhésion. Sa partie
interne, protégée par authentification, donne accès à dix-huit modules de gestion dont
l'étendue dépend du rôle effectif de l'utilisateur connecté.

Le système reconnaît neuf rôles applicatifs et cinq statuts d'appartenance, articulés par
trois types de mandats de gouvernance. Il gère 31 entités persistantes réparties en 29
tables, dont quatre tables d'historique et une table d'audit. L'ensemble représente
environ 9~900 lignes de Java côté serveur et 13~000 lignes de JavaScript côté client.

\section{Méthodologie de travail}

Le projet a été conduit de manière incrémentale, chaque incrément étant validé avant le
suivant. Quatre phases se sont succédé.

\begin{enumerate}
  \item \textbf{Analyse et modélisation.} Lecture du cahier des charges, identification
        des acteurs et des règles métier, établissement du modèle de données. C'est à ce
        stade qu'a été prise la décision structurante de faire du rôle une donnée dérivée
        plutôt qu'une donnée saisie.
  \item \textbf{Construction du socle.} Schéma relationnel versionné par Flyway, entités
        JPA, couche de sécurité, puis modules métier ajoutés un à un selon un découpage
        constant contrôleur~/~service~/~repository.
  \item \textbf{Réalisation de l'interface.} Mise en place d'un système de composants
        réutilisables, puis des 21 pages, chacune connectée à des endpoints réels --- le
        cahier des charges proscrivant explicitement les maquettes et les données fictives.
  \item \textbf{Validation.} Audit fonctionnel de bout en bout sur les 21 routes et les
        neuf rôles, correction des défauts constatés, puis écriture de tests automatisés
        verrouillant les règles métier les plus sensibles.
\end{enumerate}

Cette dernière phase a joué un rôle particulier. Elle a mis au jour des défauts qu'aucune
lecture de code n'avait révélés : une dérive de rôle sur des comptes créés avant
l'introduction du service de dérivation, des réponses d'erreur mal formées sur le chemin
d'authentification, et des commandes visibles par un vice-directeur alors que le serveur
les lui refusait. Ces corrections sont documentées au chapitre~\ref{ch:tests}.

\section{Organisation du rapport}

Le rapport suit la progression du projet. Le chapitre~\ref{ch:besoins} analyse les besoins
et formalise les acteurs et les exigences. Le chapitre~\ref{ch:conception} présente
l'architecture et la modélisation UML. Le chapitre~\ref{ch:techno} justifie les choix
technologiques, en particulier le rattachement de Spring~Boot à l'écosystème JEE. Les
chapitres~\ref{ch:backend} et~\ref{ch:frontend} détaillent la réalisation du serveur puis
de l'interface. Le chapitre~\ref{ch:securite} traite de la sécurité et de la traçabilité.
Le chapitre~\ref{ch:tests} expose la stratégie de validation et ses résultats. Le
chapitre~\ref{ch:resultats} confronte le réalisé au cahier des charges et énonce les
limites, avant le bilan du chapitre~\ref{ch:conclusion}.
